\documentclass[../midgard.tex]{subfiles}
\graphicspath{{\subfix{../images/}}}
\begin{document}

\section{Protocol security}
\label{h:protocol-security}

\subsection{L1 settlement, finality, and rollback}
\label{h:malicious-majority-stake-attack}

Midgard's canonical protocol state is evolved by Cardano L1 transactions.
An adversary that can violate Cardano's chain-selection, common-prefix, or availability guarantees can therefore affect Midgard commitments and disputes.
Midgard does not turn Cardano's probabilistic settlement into absolute or instantaneous finality.

An L2 transaction normally has two relevant L1 observations: its block header is first committed to the state queue, and the mature header is later merged into confirmed state.
The maturity delay provides time for DA retrieval and challenges, but elapsed wall-clock time alone does not make either L1 transaction irreversible.
Clients and operators must identify observations by Cardano chain point, apply a documented confirmation-depth policy, handle rollbacks of commits, proofs, removals, and merges, and quarantine provider disagreement.

A rollback of a merge may restore a header to the queue; a rollback of the earlier commit may remove the header entirely.
Neither case is harmless by definition: offchain ledgers, DA records, user receipts, settlements, payouts, and pending-operation journals must reconcile to the surviving L1 chain before production resumes.
Recommitment or remerge may be possible, but the runtime must prove that it cannot create a competing history, double-apply an event, or acknowledge a reverted result as final.

The deployment's maturity and client-confirmation policies must be derived from Cardano's current stability assumptions together with worst-case DA fetch, proof construction, multi-transaction proof execution, congestion, retry, and rollback margins.
Fixed block counts or durations in examples are not permanent security constants.

\subsection{Operator abuse of power}
\label{h:operator-abuse-of-power}

Each operator has the exclusive privilege to commit blocks to the state queue and resolve settlements during its assigned shifts (see \cref{h:time-model}), whose duration is controlled by \code{shift\_duration} (see \cref{h:protocol-parameters}).
The current operator chooses whether and when to commit and which admissible events and transactions to include.

Deposits and withdrawal orders receive L1 inclusion times.
A covering block that omits an authenticated event must be challengeable, so a continuing sequence of valid blocks cannot censor those events.
Users can also escalate an ignored L2 request through an L1 transaction order; this creates an inclusion obligation rather than an unconditional confirmation guarantee.
These protections depend on complete, correctly bound omission proofs and available L1 event data.

If all operators stop committing, inclusion-time rules do not provide liveness.
The target escape hatch temporarily reduces the bond required for replacement operators, but it is not a direct force-exit.
A production deployment must either implement and exercise that path or explicitly document the operator-liveness and fund-recovery limitation.

\subsection{Proof, DA, and economic assumptions}

Optimistic safety requires more than a slashable bond.
For every enabled state transition, an independent party must be able to retrieve the committed data, reconstruct the exact versioned rules, detect a violation, construct an admissible proof, and complete all required L1 transactions before the header becomes mergeable.
Missing proof coverage, ambiguous root binding, unavailable data, or a challenge window shorter than the worst-case proof deadline invalidates the corresponding safety claim.

Bond, slash, and reward values must be compared with maximum extractable value, detection probability, watcher capital and operating costs, proof failure risk, and correlated operator/DA failures.
Zero-valued or demonstration parameters provide no economic deterrence and must not be used to make production security claims.

\subsection{Replay resistance}
\label{h:replay-attack}

Midgard protocol actions consume Cardano UTxOs, so exact replay of a transaction on the same surviving L1 history is rejected by Cardano's double-spend rules.
Protocol-level replay resistance additionally requires deployment-domain separation, strict transaction and event identifiers, exact-once state transitions, and rollback-safe journals.

Hashed block contents are not validated merely because a header is on L1.
Duplicate transactions or events, reused roots, and mismatched proof witnesses must be rejected by the versioned ledger rules or covered by an onchain fraud proof.
Deposit absorption and withdrawal payout must also consume the unique authenticated L1 event or settlement capability so value cannot be released twice.

\end{document}
